\documentclass[11pt]{article}
\usepackage[T1]{fontenc}
\usepackage{textcomp}
\usepackage[margin=0.75in]{geometry}
\usepackage{amsmath,amssymb,amsthm}
\usepackage{array,booktabs}
\usepackage{hyperref}
\setlength{\parindent}{0pt}
\newtheorem{theorem}{Theorem}
\theoremstyle{definition}
\newtheorem{definition}{Definition}
\title{Flow Proof Artifact: List-len-rev}
\author{Generated by \texttt{flow doc proof}}
\date{Flow Proof Book}
\begin{document}
\maketitle
Reversing preserves list length.\\[0.5em]
\textbf{Source.} church — https://en.wikipedia.org/wiki/Length\_of\_a\_list\\[0.5em]
\bigskip
\noindent{\large \textbf{Derived fact 1.} \textit{len(rev(xs)) = len(xs)} \hfill List \cdot length \cdot reverse preserves length}\par
\smallskip\noindent\textit{Coordinate: List \cdot length \cdot reverse preserves length}\par
\smallskip\noindent\textit{Source: church/induction}\par
\smallskip\noindent\textit{Built on: reversing empty yields empty, for reverse on List, length adds under append, for length on List, reverse distributes over append, for reverse on List}\par
\medskip
\noindent\textbf{Goal.} len(rev(xs)) = len(xs)\par
\noindent$\displaystyle \forall xs \in List\quad len(rev(xs)) = len(xs)$\par
\medskip
\noindent
\begin{tabular}{@{} >{\raggedright\arraybackslash}p{0.06\textwidth} >{\raggedright\arraybackslash}p{0.47\textwidth} >{\raggedleft\arraybackslash}p{0.41\textwidth} @{}}
\textbf{\#} & \textbf{Proof} & \textbf{Mathematics} \\
\midrule
\textbf{1.} & We prove that reverse preserves length for length on List. & \\[0.45em]
\textbf{2.} & We proceed by induction on xs: first the base case, then the inductive step. & \\[0.45em]
\textbf{3.} & Consider the base case in which xs  equals  nil. & \\[0.45em]
\textbf{4.} & We invoke the derived fact governing reverse on List: reversing empty yields empty, for reverse on List. & \\[0.45em]
\textbf{5.} & We invoke the derived fact governing length on List: the empty list has length zero, for length on List. & \\[0.45em]
\textbf{6.} & From step 3, step 4, and step 5, we can deduce that len(rev(xs)) equals len(xs). This establishes the base case (see step 3, step 4, and step 5). Hence proven. & $\displaystyle len(rev(xs)) = len(xs)$ \\[0.45em]
\textbf{7.} & For the inductive step, suppose the claim holds for all smaller values. & \\[0.45em]
\textbf{8.} & Under the supposition in step 7, let rest = tail(xs). & \\[0.45em]
\textbf{9.} & Under the supposition in step 7, we cross the inductive boundary: assume the claim holds for rest (the induction hypothesis). & $\displaystyle len(rev(rest)) = len(rest)$ \\[0.45em]
\textbf{10.} & We invoke the derived fact governing reverse on List: reverse distributes over append, for reverse on List (instantiated for rest, cons(head(xs), nil)). & \\[0.45em]
\textbf{11.} & From step 7, step 8, step 9, and step 10, this implies len(rev(xs)) equals len(xs). Hence proven. & $\displaystyle len(rev(xs)) = len(xs)$ \\[0.45em]
\bottomrule
\end{tabular}
\label{thm:List-length-reverse_preserves_length}
\par\medskip\hrule
\end{document}
